% 第 18 讲 范数群、Kummer 理论与投射有限上同调（译自 src/lecture18.txt）
\section*{第 18 讲\quad 范数群、Kummer 理论与投射有限上同调}

上一次我们证明了消没定理，并看到它蕴涵：对每个有限 Galois $G$-扩张 $L/K$，有 $(L^\times)_{tG}\simeq \mathbb Z_{tG}[-2]$；取零次上同调即得 $K^\times/NL^\times\simeq G^{\mathrm{ab}}$。注意后者不可能平凡，因为 $G$ 必须是可解群。然而在第一讲中，我们表述了另一个定理：
\[
\Gal(\bar K/K)^{\mathrm{ab}}:=\varprojlim_{L/K}\Gal(L/K)^{\mathrm{ab}}\simeq \widehat{K^\times},
\]
其中逆向极限跑遍有限 Galois 扩张 $L/K$。回忆
\[
\widehat{K^\times}:=\varprojlim_{[K^\times:\Gamma]<\infty}K^\times/\Gamma
\]
是 $K^\times$ 的投射有限完备化，其中 $\Gamma$ 是 $K^\times$ 的有限指数闭子群（这是对拓扑群定义投射有限完备化的唯一合理方式）。于是我们想证明
\[
\varprojlim_{L/K}K^\times/NL^\times\simeq \varprojlim_{[K^\times:\Gamma]<\infty}K^\times/\Gamma,
\]
其中 $L$ 与 $\Gamma$ 如上。

\textbf{定义 18.1.} $K^\times$ 的子群 $\Gamma$ 称为范数群（或范数子群），若对某个有限扩张 $L/K$ 有 $\Gamma=NL^\times$。

\textbf{定理 18.2（存在定理）.} $K^\times$ 的子群 $\Gamma$ 是范数群，当且仅当 $\Gamma$ 是闭的且指数有限。

这显然足以证明上述局部类域论的陈述。

\textbf{注 18.3.} 局部类域论的一个推论是：若 $L/K$ 是 $G$-Galois 扩张，而 $L/L_0/K$ 是 $K$ 在 $L$ 中的极大阿贝尔子扩张，则 $NL^\times=NL_0^\times$。这是因为
\[
K^\times/NL^\times\simeq G^{\mathrm{ab}}\simeq K^\times/NL_0^\times.
\]

我们将在 $\operatorname{char}(K)=0$ 的情形证明存在定理，尽管它在其他情形也成立（只是证明更复杂）。

\textbf{引理 18.4.} 若 $\Gamma\subseteq K^\times$ 是范数群，则 $\Gamma$ 闭且指数有限。

\textit{证明.} 设 $L/K$ 是 $n$ 次扩张，使 $\Gamma=NL^\times$。则 $\Gamma\supseteq N_{L/K}K^\times=(K^\times)^n$，而我们已经看到后者是有限指数的闭子群（因为它包含 $1+\mathfrak p_K^N$，其中 $N$ 充分大），故 $\Gamma$ 亦然。注意若 $\operatorname{char}(K)>0$，则 $(K^\times)^n$ 在 $K^\times$ 中的指数实际上是无限的！\hfill$\square$

于是存在定理的内容是：对一切 $n$，$\pi^{n\mathbb Z}(1+\mathfrak p_K^n)$ 都是范数子群；我们已经证明范数子群「不太小」，现在需要证明它们可以「足够小」。

\textbf{引理 18.5.} 若 $\Gamma'$ 是 $K^\times$ 的子群，且对某个范数子群 $\Gamma$ 有 $K^\times\supseteq\Gamma'\supseteq\Gamma$，则 $\Gamma'$ 也是范数子群。

\textit{证明.} 设 $L/K$ 是有限扩张，使 $\Gamma=NL^\times$。如前，不妨设 $L/K$ 阿贝尔。则由局部类域论，
\[
\Gamma'/\Gamma\subseteq K^\times/NL^\times\simeq\Gal(L/K)
\]
是正规子群，因为按假设 $\Gal(L/K)$ 阿贝尔。于是存在中间扩张 $K'/K$（$K\subseteq K'\subseteq L$），使 $\Gamma'/\Gamma=\Gal(L/K')$，并且
\[
K^\times/N(K')^\times=\Gal(K'/K)=\Gal(L/K)/\Gal(L/K')=(K^\times/NL^\times)/(\Gamma'/\Gamma)=K^\times/\Gamma'
\]
典范地成立。于是 $\Gamma'=N(K')^\times$，即所欲证。

注意我们隐含地使用了如下事实：按我们对局部类域论的明确构造，下述图表交换（对阿贝尔扩张 $L/K$）：
\[
\begin{CD}
\Gal(L/K) @>\sim>> K^\times/NL^\times\\
@V{\alpha}VV @V{}VV\\
\Gal(K'/K) @>\sim>> K^\times/N(K')^\times.
\end{CD}
\]
由于 $\Gamma'/\Gamma=\ker(\alpha)$ 在 $K^\times$ 中的原像既是 $\Gamma'$ 又是 $N(K')^\times$，我们再次得到 $\Gamma'=N(K')^\times$。\hfill$\square$

现在作一段离题讨论：在第二讲中我们说过
\[
K^\times/(K^\times)^2\simeq \Gal^{\mathrm{ab}}(\bar K/K)/2\simeq \operatorname{Hom}\big(K^\times/(K^\times)^2,\mathbb Z/2\mathbb Z\big),
\]
假设 $\operatorname{char}(K)=0$（特别地不为 $2$），其中第一个同构经由局部类域论。也就是说，$K^\times/(K^\times)^2$ 是自对偶的。现在问：如何把它推广到 $n=2$ 之外？答案是使用 Kummer 理论。

回忆：假设 $n\nmid \operatorname{char}(K)$，且 $n$ 次单位根群 $\mu_n\subseteq K^\times$ 的阶为 $n$，则
\[
K^\times/(K^\times)^n\simeq \operatorname{Hom}_{cts}\big(\Gal(\bar K/K),\mu_n\big),
\]
其中右端是群同态。要点是：若 $K$ 还是局部域，我们便期望
\begin{equation}
K^\times/(K^\times)^n\simeq \operatorname{Hom}\big(K^\times/(K^\times)^n,\mu_n\big). \tag{18.1}
\end{equation}

事实上，我们有如下映射：
\begin{align*}
K^\times/(K^\times)^n &= \operatorname{Hom}_{cts}\big(\Gal(\bar K/K),\mu_n\big)\\
&= \operatorname{Hom}_{cts}\big(\Gal^{\mathrm{ab}}(\bar K/K),\mu_n\big)\\
&= \operatorname{Hom}_{cts}\Big(\varprojlim_{L/K}K^\times/NL^\times,\mu_n\Big)\\
&= \varinjlim_{L/K}\operatorname{Hom}\big(K^\times/NL^\times,\mu_n\big)\\
&\hookrightarrow \operatorname{Hom}_{cts}\big(K^\times,\mu_n\big)\\
&= \operatorname{Hom}\big(K^\times/(K^\times)^n,\mu_n\big),
\end{align*}
其中第二个等号是因为所有这类映射都必须穿过 $\Gal(\bar K/K)$ 的阿贝尔化（因为 $\mu_n$ 阿贝尔），第三个由局部类域论，第四个由对偶性。注意这些逆向极限跑遍有限扩张 $L/K$；而「连续」（当定义域有限时并不需要）在这里指映射杀掉某个紧开子群，这就解释了上面的单射。我们想证明这个映射也是同构。注意 $K^\times/(K^\times)^n$ 是有限阿贝尔群且为 $n$-挠；因此只需证明两边阶相同。

\textbf{断言 18.6.} 设 $A$ 是 $n$-挠有限阿贝尔群。则
\[
\#A=\#\operatorname{Hom}(A,\mathbb Z/n\mathbb Z).
\]

\textit{证明.} $A$ 是若干 $\mathbb Z/d\mathbb Z$（$d\mid n$）的直和，故可归约到 $A=\mathbb Z/d\mathbb Z$ 的情形（一般情形下直和与 $\operatorname{Hom}$ 交换）。于是
\[
\operatorname{Hom}(\mathbb Z/d\mathbb Z,\mathbb Z/n\mathbb Z)=(\mathbb Z/n\mathbb Z)[d],
\]
其阶为 $d=\#A$，如所欲证。\hfill$\square$

这表明 (18.1) 是典范同构（不过仅凭断言的一般陈述也足以说明它是同构）。在 $n=2$ 的情形，容易看出这正是希尔伯特符号。

\textbf{推论 18.7.} 若 $\mu_n\subseteq K$，则 $(K^\times)^n$ 是范数子群。

\textit{证明.} 若把我们 Kummer 理论的「图景」对偶化，便得到如下交换图，连同自然的连续映射 $\alpha:K^\times\to\varprojlim_{L/K}K^\times/NL^\times$：
\[
\begin{CD}
\Gal(\bar K) @>{cts}>> \operatorname{Hom}\big(K^\times/(K^\times)^n,\mu_n\big)\\
@V{\text{投影}}VV @A{\beta_{cts}}AA\\
\varprojlim_{L/K}K^\times/NL^\times @>{\sim}>> \Gal^{\mathrm{ab}}(\bar K)
\end{CD}
\]
% [原图在 PDF 中为「Gal(K) --cts--> Hom(K^×/(K^×)^n,μn)」与「K^× --α_cts--> lim K^×/NL^× = Gal^ab(K)」两行、右侧竖直箭头 β_cts（向上）、中间竖直箭头 Gal(K)→Gal^ab(K)，已按原 PDF 第 82 页图式复原]
注意 $\alpha$ 连续，因为逆向极限中的开子群都是范数子群，从而它在 $K^\times$ 中的原像是有限指数开子群。正如我们刚刚看到的，$\ker(\beta\circ\alpha)=(K^\times)^n$，它在逆向极限中是开的（即它是某个 $L/K$ 的 $K^\times/NL^\times$ 的某个子集沿典范投影的完全原像，因为这些映射连续）。于是由引理 18.5，$(K^\times)^n$ 是范数子群。\hfill$\square$

注意上面的映射 $\beta$ 是满射，因为它实现为 $\Gal^{\mathrm{ab}}(\bar K/K)$ 模去 $n$ 次幂。

\textbf{注 18.8.}「先验地」（即如果我们忘掉每个群的阶），这个复合的核可能比 $(K^\times)^n$ 更大。通过论证二者相等，我们便得到了一个「小」的范数子群。

\textit{证明（存在定理）.} 设 $K$ 是一般的特征 $0$ 局部域。令 $L:=K(\zeta_n)$，考虑扩张 $L/K$，其中 $\zeta_n$ 表示本原 $n$ 次单位根的集合。由于由推论 18.7，$(L^\times)^n$ 是 $L^\times$ 中的范数子群，故 $N(L^\times)^n=N((L^\times)^n)\subseteq K^\times$ 是 $K^\times$ 中的范数子群。但 $N(L^\times)^n\subseteq (K^\times)^n\subseteq K^\times$，于是引理 18.5 表明 $(K^\times)^n$ 是 $K^\times$ 中的范数子群。

现在注意，对一切 $N$，存在某个 $n$ 使得
\[
(\mathcal O_K^\times)^n=(K^\times)^n\cap\mathcal O_K^\times\subseteq 1+\mathfrak p_K^N.
\]
事实上，$(\mathcal O_K^\times)^{q-1}\subseteq 1+\mathfrak p_K$，其中 $q=\#\mathcal O_K/\mathfrak p_K$（因为模 $\mathfrak p_K$ 的剩余类取 $q-1$ 次幂必为 $1$）。于是对充分大的 $v(n)$ 有 $(\mathcal O_K^\times)^{(q-1)n}\subseteq 1+\mathfrak p_K^N$，因为一般地 $(1+x)^n=1+nx+\cdots$（省略号表示高阶项），而当 $v(n)\gg 0$ 时除 $1$ 以外的所有项都落在 $\mathfrak p_K^N$ 中。

至于「$\mathbb Z$-方向」上的有限指数子群，即限制到 $\pi^N$ 的倍数上的情形，只需把 $n$ 换成 $nN$，使得只有赋值被 $N$ 整除的元素被实现。于是 $K^\times$ 的每个有限指数开子群都包含某个 $(K^\times)^n$，而后者如上所示是范数子群，故由引理 18.5 它本身也是范数子群。\hfill$\square$

现在让我们快速重访 Kummer 理论；我们将说明，它实际上关于群上同调讲了一件非常一般的事情。设 $G$ 是投射有限群，即 $G=\varprojlim_i G_i$，其中 $G_i$ 是有限群。

\textbf{定义 18.9.} 称 $G$-模 $M$ 是光滑的，若对一切 $x\in M$，存在有限指数开子群 $K\subseteq G$ 使得 $K\cdot x=x$。

\textbf{例 18.10.} 若 $G:=\Gal(\bar K/K)$，则 $G$ 同时作用在 $\bar K$ 与 $\bar K^\times$ 上，二者都是光滑 $G$-模。这是因为这两个 $G$-模中的每个元素都落在某个有限扩张 $L/K$ 里，从而被 $\Gal(\bar K/L)$ 固定，而后者按定义是有限指数开子群。

光滑性使我们能从往往非常复杂的投射有限群归约到有限群的情形。我们现在必须为投射有限群定义一个群上同调的概念，因为我们原先的表述只针对有限群。

\textbf{定义 18.11.} 设 $X$ 是下有界的光滑 $G$-模复形。则
\[
X^{hG}:=\varinjlim_i (X^{K_i})^{hG/K_i},
\]
其中 $K_i:=\ker(G\to G_i)$，而 $X^{K_i}$ 表示被 $K_i$（朴素地）固定的向量。容易看出这构成正向系统。注意 $G_i$ 不作用在 $X$ 上（它只是 $G$ 的商），但它作用在被 $K_i$ 固定的向量上。这些 $K_i$ 是 $G$ 的紧开子群，大小递减。用 $K_i$ 取「朴素不变量」是令人担心的，因为它不保持拟同构；但实际上我们有：

\textbf{断言 18.12.} 若 $X$ 无环，则 $X^{hG}$ 也无环。

证明略去；不过要注意，$X$ 下有界这一点很重要。我们有 Hilbert 定理 90 的如下「无限版本」：

\textbf{命题 18.13.} 若 $L/K$ 是（可能无限的）$G$-Galois 扩张，则
\[
H^1(G,L^\times):=H^1\big((L^\times)^{hG}\big)=0.
\]

\textit{证明.} 写 $L=\bigcup_i L_i$，其中每个 $L_i$ 是 $K$ 的有限 $G_i$-Galois 扩张。则按定义，
\[
H^1(G,L^\times)=\varinjlim_i H^1(G_i,L_i^\times)=0
\]
由 Hilbert 定理 90。\hfill$\square$
% [原文作 $H^1(G_i,K_i^\times)$，但按 $L=\bigcup_i L_i$（各 $L_i$ 为 $K$ 的有限 $G_i$-Galois 扩张）及定义 18.11（其中 $K_i=\ker(G\to G_i)$ 另有含义）应为 $L_i^\times$；已订正]

\textbf{推论 18.14.} 设 $G:=\Gal(\bar K/K)$，$n$ 与 $\operatorname{char}(K)$ 互素。若 $\mu_n\subseteq K$，则
\[
K^\times/(K^\times)^n\simeq \operatorname{Hom}_{cts}(G,\mu_n).
\]

\textit{证明.} 我们有光滑 $G$-模的短正合列
\[
0\to \mu_n\to \bar K^\times\xrightarrow{\ x\mapsto x^n\ }\bar K^\times\to 0.
\]
上同调的长正合列给出
\[
K^\times\xrightarrow{\ x\mapsto x^n\ }K^\times\to H^1(G,\mu_n)\to H^1(G,\bar K^\times),
\]
其中两端的 $K^\times$ 即 $H^0(G,\bar K^\times)$，而最右端由 Hilbert 定理 90（命题 18.13）为零。于是 $K^\times/(K^\times)^n\simeq H^1(G,\mu_n)$。由于 $\mu_n\subseteq K$（如 Kummer 理论的设定），它被 $G$ 固定；正如我们通过上链看到的，对平凡群作用有 $H^1(G,\mu_n)=\operatorname{Hom}_{cts}(G,\mu_n)$，这就给出所要结果。\hfill$\square$

于是我们确实可以从抽象群上同调与 Hilbert 定理 90 非常简单地推出 Kummer 理论。
